\chapter{Superviser l'exploitation}

\section{Alertes, interventions et maintenance}

Ces trois modules sont liés mais ne représentent pas la même étape du travail.

\begin{longtable}{|L{3.1cm}|L{4.7cm}|L{4.5cm}|L{3.2cm}|}
  \hline
  \rowcolor{ctmint}
  \textbf{Module} & \textbf{Objet} & \textbf{Exemple} & \textbf{Responsable
  principal} \\ \hline
  \endfirsthead
  \hline
  \rowcolor{ctmint}
  \textbf{Module} & \textbf{Objet} & \textbf{Exemple} & \textbf{Responsable
  principal} \\ \hline
  \endhead
  Alerte &
  Signaler automatiquement ou manuellement un événement à analyser. &
  Perte de connexion, défaut connecteur, température anormale. &
  Opérateur. \\ \hline
  Intervention &
  Organiser l'action humaine qui traite un incident ou une tâche. &
  Diagnostiquer et remplacer un composant. &
  Opérateur puis technicien assigné. \\ \hline
  Maintenance &
  Planifier et tracer les opérations préventives ou correctives. &
  Inspection périodique ou entretien après panne. &
  Administrateur ou opérateur. \\ \hline
\end{longtable}

\subsection{Workflow recommandé}

\begin{procedurebox}{Traiter un incident opérationnel}
  \begin{enumerate}[leftmargin=*]
    \item Examiner l'alerte, sa gravité, la station et l'événement source.
    \item Accuser réception afin d'éviter un traitement parallèle inutile.
    \item Résoudre à distance si l'action est sûre et suffisante.
    \item Sinon, créer une intervention, définir la priorité et l'échéance,
    puis l'assigner à un technicien.
    \item Le technicien effectue le diagnostic, joint les preuves avant/après et
    soumet son rapport.
    \item Valider le résultat, clôturer l'intervention puis résoudre l'alerte.
    \item Créer une maintenance planifiée si une action de suivi est nécessaire.
  \end{enumerate}
\end{procedurebox}

\section{Tarifs}

Un tarif définit les règles utilisées pour estimer et facturer une recharge :
prix de l'énergie, frais fixes, éventuels frais de durée et conditions
d'application. Les modifications doivent être datées et ne doivent pas changer
rétroactivement le montant d'une session déjà facturée.

\begin{procedurebox}{Publier un tarif}
  \begin{enumerate}[leftmargin=*]
    \item Ouvrir \emph{Tariffs \& plans}.
    \item Créer un brouillon avec une désignation explicite.
    \item Renseigner les composantes tarifaires et leur unité.
    \item Définir la période et les stations concernées.
    \item Vérifier un exemple de calcul.
    \item Activer le tarif après validation interne.
  \end{enumerate}
\end{procedurebox}

\section{Sessions et paiements}

La session contient les mesures reçues pendant la recharge, le tarif appliqué
et le montant calculé. Le paiement client est géré séparément de la facture
commerciale de l'organisation.

Dans le MVP, l'adaptateur de paiement externe est simulé. Les états
d'autorisation, de confirmation et d'échec permettent de tester le workflow,
mais ne représentent aucun transfert bancaire réel.

\section{Rapports et exports}

L'espace d'analyse de l'administrateur couvre les résultats commerciaux,
l'utilisation du réseau, les sessions, les alertes, les interventions et la
maintenance de son organisation.

Les boutons d'export proposent les formats CSV, JSON et PDF lorsque ces formats
sont adaptés au contenu. Le PDF est destiné à la lecture et à l'archivage ; le
CSV et le JSON servent plutôt à l'analyse ou à l'intégration.

\subsection{Échange de rapports}

Les employés et l'administrateur peuvent composer et échanger des rapports
dans le périmètre de l'organisation. Un rapport doit inclure un objet, une
période, les entités concernées, une conclusion et, si nécessaire, des pièces
jointes.

\securitynote{Vérifier les destinataires avant l'envoi. Une pièce jointe peut
contenir une photo d'équipement, une position ou une information opérationnelle
sensible.}

\manualscreenshot{admin-reports.png}
  {Analyses et rapports de l'administrateur}
  {fig:admin-reports}
